% To be reformatted in the TAC style (tac.cls) upon acceptance.
\documentclass[11pt]{amsart}

\usepackage{amsmath,amssymb}
\usepackage[hidelinks]{hyperref}

\newtheorem{thm}{Theorem}[section]
\newtheorem{lem}[thm]{Lemma}
\newtheorem{cor}[thm]{Corollary}
\newtheorem{prop}[thm]{Proposition}
\theoremstyle{definition}
\newtheorem{defi}[thm]{Definition}
\newtheorem{rem}[thm]{Remark}

\newcommand{\I}{\mathbb{I}}
\newcommand{\DM}[1]{\mathrm{DM}(#1)}
\newcommand{\Dbox}{\square_{\mathrm{DM}}}
\newcommand{\aconv}{\equiv}
\newcommand{\nconv}{\not\equiv}
\newcommand{\propeq}{\simeq}
\newcommand{\refl}{\mathsf{refl}}
\newcommand{\Path}{\mathsf{Path}}
\newcommand{\PathP}{\mathsf{PathP}}
\newcommand{\U}{\mathcal{U}}
\newcommand{\hcomp}{\mathsf{hcomp}}
\newcommand{\hfillop}{\mathsf{hfill}}
\newcommand{\act}[1]{A_{#1}}
\newcommand{\itp}[2]{#1 \mathbin{\ast_k} #2}

\begin{document}

\title[De Morgan convexity]{De Morgan convexity and the
  strict--weak comparison\\ in cubical type theory}

\author{Ryota Shibusawa}
\address{Daiichi Institute of Technology, Japan}
\email{r-shibusawa@daiichi-koudai.ac.jp}

\subjclass[2020]{18N45, 03B38, 18M99}
\keywords{cubical type theory, connections, De Morgan algebra,
  definitional equality, homotopy, Kan filling, excluded middle}

\begin{abstract}
  A companion paper proved that in cubical type theory the
  \emph{strict fillers} of a boundary --- the fillers definable by
  pure reparametrization, without Kan operations --- need not be
  unique.  This paper shows that the non-uniqueness resolves
  \emph{inside} the strict layer: the De Morgan convex combination
  \[
    \chi(k) \;=\; (\varphi \wedge \neg k) \vee (k \wedge \varphi')
    \vee (\varphi \wedge \varphi')
  \]
  yields a strict homotopy rel boundary between any two
  same-boundary fillers of one cell, and, iterated, makes the fiber
  of strict fillers over a boundary strictly contractible --- the
  exact strict parallel of the uniqueness of Kan fillers up to
  homotopy.  The third term is essential: the naive combination
  $(\varphi \wedge \neg k) \vee (k \wedge \varphi')$ fails
  \emph{precisely} by the absence of excluded middle in the free De
  Morgan algebra, and the failure is machine-checked as a typing
  error.  Fillers originating from different cells admit no direct
  strict homotopy, but connect by deformation onto the shared
  subcomplex (a machine-checked zig-zag).  The strict--weak
  comparison then closes a circle opened by the first paper of this
  series: the constant-tube composition equation, which its no-go
  theorem shows can never be definitional, holds
  \emph{propositionally and uniformly} via the canonical filler ---
  what cannot be definitional is exactly what $\hfillop$ supplies
  propositionally --- while genuinely weak composites open new
  propositional classes (winding).  All positive and negative
  instances are machine-checked in a self-contained proof-assistant
  kernel, and every statement is a decidable, assumption-free fact
  about the algorithmic equality of a minimal fragment.
\end{abstract}

\maketitle

\section{Introduction}
\label{sec:intro}

This is the fourth paper in a series on the definitional (strict)
structure of cubical type theory \cite{CCHM}, written, like its
predecessors, to be read independently.  The first paper
\cite{Const} developed the metatheoretic platform --- in particular
a no-go theorem: extending definitional equality with
\emph{constant-tube collapse} for homogeneous composition (the
equation we call $(\dagger)$ below) is impossible in principle ---
and the decidability of the algorithmic equality on a minimal
fragment.  The second \cite{Sphere} characterized the strict layer
of generic loops as a wedge of De Morgan spheres.  The third
\cite{Realization} extended the characterization to cells with
generic boundaries --- the strict layer of a cellular context
realizes its cube presheaf on the nose --- and isolated a phenomenon
absent from the Kan world: \emph{strict fillers of a common boundary
need not be unique}.

The present paper is about what that non-uniqueness means, and more
generally about how the strict layer sits inside the full theory.
Its results:

\begin{enumerate}
\item \emph{De Morgan interpolation}
  (Section~\ref{sec:interp}).  For interval formulas
  $\varphi, \varphi'$ whose restrictions to every face of the cube
  agree, the \emph{De Morgan convex combination}
  \[
    \chi(k) := (\varphi \wedge \neg k) \vee (k \wedge \varphi')
    \vee (\varphi \wedge \varphi')
  \]
  satisfies $\chi(0) = \varphi$, $\chi(1) = \varphi'$ and is
  constant in $k$ on every face.  Consequently any two strict
  fillers of one boundary arising from the same cell are connected
  by a \emph{strict} homotopy rel boundary --- no Kan operation is
  involved (machine-checked on the non-uniqueness pair of
  \cite{Realization}).
\item \emph{Necessity of the third term}
  (Section~\ref{sec:interp}).  The naive combination
  $(\varphi \wedge \neg k) \vee (k \wedge \varphi')$ restricts on a
  face with common value $\rho$ to $\rho \wedge (k \vee \neg k)$,
  which is \emph{not} $\rho$: the free De Morgan algebra has no
  excluded middle.  The would-be homotopy fails to typecheck
  (machine-checked failure).  The correction term
  $\varphi \wedge \varphi'$ is the exact price of the missing
  excluded middle.
\item \emph{Strict contractibility of fiber towers}
  (Section~\ref{sec:contract}).  Iterating the interpolation ---
  two homotopies between the same fillers again share a boundary
  --- the entire tower of strict fillers, homotopies, homotopies
  between homotopies, \dots\ over a fixed boundary is strictly
  contractible.  Strict fillers are therefore \emph{unique up to
  strict homotopy}: the exact parallel, inside the strict layer, of
  the Kan world's uniqueness of fillers up to homotopy.
\item \emph{Cross-cell connection through the shared subcomplex}
  (Section~\ref{sec:cross}).  Fillers of one boundary originating
  from \emph{different} cells admit no direct strict homotopy (a
  strict cube is single-cell material), yet connect by strict
  deformations onto the shared boundary material --- a
  machine-checked zig-zag
  $S_1 \rightsquigarrow S_m$, $S_2 \rightsquigarrow S_m$.
\item \emph{The comparison closed}
  (Section~\ref{sec:comparison}).  The constant-tube equation
  $(\dagger)$, provably never definitional \cite{Const}, holds
  propositionally and uniformly: the canonical filler is a path rel
  boundary from the base to the composite (machine-checked).
  Conversely the weak layer genuinely extends the propositional
  classification: the composite $\mathsf{loop} \cdot \mathsf{loop}$
  on the circle has winding $+2$, separating it propositionally
  from the basic strict cells (winding $0, \pm 1$,
  machine-computed).
\end{enumerate}

The emerging picture is summarized in
Table~\ref{tab:comparison}: over a common boundary, strict cells
are unique up to strict homotopy; constant-tube composites are
definitionally new but propositionally strict; genuine composites
are new outright.  The interpolation operator equips the strict
fibers with a canonical convex structure --- we call it \emph{De
Morgan convexity} --- which supplies the strict layer with its own
homotopy theory, definable entirely in the connection algebra.

\begin{table}
  \centering
  \footnotesize
  \begin{tabular}{lll}
    cells over a common boundary & definitional & propositional \\
    \hline
    strict (reparametrizations) &
      unique up to strict homotopy &
      likewise \\
    constant-tube composites &
      separated ($(\dagger)$ no-go) &
      strict, via the filler \\
    genuine composites &
      separated &
      new classes (winding) \\
  \end{tabular}
  \medskip
  \caption{The strict--weak comparison over a common boundary.}
  \label{tab:comparison}
\end{table}

\subparagraph*{Honest positioning.}  Each ingredient is elementary:
the interpolation lemma is three absorption computations, and the
propositional half of the comparison is the standard filler.  The
contribution is the assembly and its exactness --- that the
interpolant is boundary-constant \emph{exactly} when restrictions
agree, that the naive interpolant fails \emph{exactly} by no-LEM,
that strict homotopy uniqueness holds \emph{despite} definitional
non-uniqueness, and that the definitional/propositional boundary
traced by the no-go theorem of \cite{Const} is exactly the
$\hfillop$-image.  We know of no prior treatment of the strict
layer's internal homotopy theory.  As throughout the series, all
instances are machine-checked and, by the decidability theorem of
\cite{Const}, decided.

\section{Preliminaries}
\label{sec:prelim}

We work in the reparametrization-plus-composition fragment of a
CCHM-style cubical type theory \cite{CCHM}: path types over a base
type, path abstraction/application, and homogeneous composition
$\hcomp$; $\DM{\vec i}$ is the free De Morgan algebra of interval
formulas (no excluded middle), with canonical antichain normal
forms.  Definitional equality $\aconv$ is the algorithmic equality
of the kernel of \cite{Const}; we import from the companions:

\begin{enumerate}
\item[(K1)] endpoint collapse of path application is normal-form
  exact, and for dependent path types returns the stored boundary
  cells \cite[K1]{Realization};
\item[(K2)] conversion of neutral spines is head- and
  argument-exact \cite[K2]{Sphere};
\item[(K3)] on the minimal fragment, evaluation and conversion
  terminate and $\aconv$ is decidable, with no assumptions
  \cite{Const};
\item[(K4)] \emph{realization} \cite{Realization}: over a cellular
  context, strict cubes are classified, injectively, by cells and
  formula tuples; in particular \emph{equal boundaries force equal
  face restrictions} for same-cell fillers.
\end{enumerate}

\subparagraph*{The running example.}  Over the generic-boundary
square computad (vertices, edges $p,q,r,s$, square $Q$)
\cite{Realization} exhibits two strict fillers of one boundary:
\[
  F_1 := \langle i \rangle \langle j \rangle\,
  Q(i \wedge \neg i)(j),
  \qquad
  F_2 := \langle i \rangle \langle j \rangle\,
  Q\bigl((i \wedge \neg i) \wedge (j \vee \neg j)\bigr)(j),
\]
non-convertible, both inhabiting
$\PathP\,(\langle i \rangle\,
\Path\,A\,(r(i \wedge \neg i))\,(s(i \wedge \neg i)))\,p\,p$:
the two formulas differ in the free algebra yet restrict equally to
all four faces.

\section{De Morgan interpolation}
\label{sec:interp}

\begin{lem}[Interpolation]\label{lem:interp}
  For $\varphi, \varphi' \in \DM{\vec i}$ let
  \[
    \chi(k) \;:=\; (\varphi \wedge \neg k) \vee
    (k \wedge \varphi') \vee (\varphi \wedge \varphi')
    \;\in\; \DM{\vec i, k}.
  \]
  Then:
  \begin{enumerate}
  \item $\chi(0/k) = \varphi$ and $\chi(1/k) = \varphi'$;
  \item for any face $f$ of the $\vec i$-cube on which
    $\varphi|_f = \varphi'|_f = \rho$:\ \ $\chi|_f = \rho$,
    constant in $k$.
  \end{enumerate}
\end{lem}
\begin{proof}
  (1) $\chi(0) = (\varphi \wedge 1) \vee 0 \vee (\varphi \wedge
  \varphi') = \varphi \vee (\varphi \wedge \varphi') = \varphi$ by
  absorption ($\varphi \wedge \varphi' \le \varphi$); symmetrically
  at $1$.  (2) $\chi|_f = (\rho \wedge \neg k) \vee (k \wedge \rho)
  \vee \rho = \rho$, again by absorption (the first two joinands are
  $\le \rho$).
\end{proof}

\begin{thm}[Strict homotopy rel boundary]\label{thm:homotopy}
  Let $F = \act{c}(\vec\varphi)$ and $F' = \act{c}(\vec\varphi')$
  be strict fillers of the same boundary arising from the same cell
  $c$.  Then
  \[
    H \;:=\; \act{c}(\vec\chi), \qquad
    \chi_t := (\varphi_t \wedge \neg k) \vee (k \wedge \varphi_t')
    \vee (\varphi_t \wedge \varphi_t'),
  \]
  with $k$ as the outer dimension, is a strict homotopy rel
  boundary from $F$ to $F'$.
\end{thm}
\begin{proof}
  By (K4), equality of the boundaries forces
  $\vec\varphi|_f = \vec\varphi'|_f$ on every face $f$.
  Lemma~\ref{lem:interp}(1) gives the $k$-endpoints;
  Lemma~\ref{lem:interp}(2), applied coordinatewise, makes every
  face of $H$ constant in $k$ and equal to the given boundary, so
  $H$ inhabits the path type over the boundary type from $F$ to
  $F'$.  (Typing uses only the boundary conditions; interior
  degeneracies of $\vec\chi$ are irrelevant.)  Machine witness on
  the running example: \texttt{mixInterpD}, with a control probe
  and with the separation $F_1 \nconv F_2$ re-checked --- the
  homotopy is not trivial.
\end{proof}

\begin{lem}[Necessity of the third term]\label{lem:necessity}
  For $k$-free $\rho \ne 0$:
  \[
    (\rho \wedge \neg k) \vee (k \wedge \rho)
    \;=\; \rho \wedge (k \vee \neg k) \;\ne\; \rho ,
  \]
  since $\rho \le k \vee \neg k$ would require every clause of
  $\rho$ to contain $k$ or $\neg k$.  Consequently the naive
  interpolant $(\varphi \wedge \neg k) \vee (k \wedge \varphi')$ is
  \emph{not} boundary-constant, and the corresponding term fails to
  typecheck as a homotopy (machine-checked failure on the running
  example).
\end{lem}

\begin{rem}
  Lemma~\ref{lem:necessity} is a small but pointed fact: in a
  Boolean interval the naive convex combination would do, and the
  third term is invisible.  The free De Morgan algebra's missing
  excluded middle --- the very feature that keeps the connection
  layer faithful \cite{Sphere} --- exacts the correction term
  $\varphi \wedge \varphi'$ as its price.  The corrected operation
  is the ``strong'' (Kleene-style) convex combination, and we will
  see it is exactly what the strict layer's homotopy theory runs
  on.
\end{rem}

\section{Strict contractibility of fiber towers}
\label{sec:contract}

Fix a cellular context, a cell $c$, and a boundary $\beta$.

\begin{defi}[Fiber tower]\label{def:tower}
  $\mathrm{Fib}^0(\beta)$ is the set of strict fillers of $\beta$
  from $c$; for $x, y \in \mathrm{Fib}^n$,
  $\mathrm{Fib}^{n+1}(x,y)$ is the set of strict homotopies rel
  boundary between them (strict cubes one dimension up whose
  $k$-faces are $x, y$ and whose remaining faces are the constant
  extensions of the lower boundary data).  The tower is
  \emph{strictly contractible} if at every level any two elements
  are connected: $\mathrm{Fib}^{n+1}(x,y) \ne \emptyset$ for all
  $x, y$.
\end{defi}

\begin{thm}[Contractibility]\label{thm:contract}
  The fiber tower of a cell over any boundary is strictly
  contractible.
\end{thm}
\begin{proof}
  Level $1$ is Theorem~\ref{thm:homotopy}.  For level $n+1$: two
  elements $H, H' \in \mathrm{Fib}^n$ are strict cubes
  $\act{c}(\vec\chi), \act{c}(\vec\chi')$ of the same $(m+n)$-cube
  boundary --- their $k$-faces are the shared lower-level elements
  and their remaining faces the shared constant extensions --- so
  by (K4) their tuples' restrictions agree on \emph{every} face,
  and Theorem~\ref{thm:homotopy} applies verbatim one dimension
  up.  Induction.
\end{proof}

\begin{rem}
  Together with the non-uniqueness theorem of \cite{Realization},
  the picture is: \emph{strict fillers are not unique, but are
  unique up to strict homotopy, coherently at all levels} --- the
  precise strict parallel of Kan filling's uniqueness up to
  homotopy, achieved without any Kan operation.  We call the
  underlying canonical structure --- the interpolation operator as
  a convex combination on fibers --- \emph{De Morgan convexity}.
\end{rem}

\section{Cross-cell fillers and the shared subcomplex}
\label{sec:cross}

Fillers of one boundary can originate from different cells: over
the shared-edge computad of \cite{Realization} (squares $Q_1, Q_2$
sharing the edge $m$), the squares
\[
  S_1 := \langle i \rangle \langle j \rangle\,
  Q_1(i \vee \neg i \vee j \vee \neg j)(j),
  \qquad
  S_2 := \langle i \rangle \langle j \rangle\,
  Q_2(i \wedge \neg i \wedge j \wedge \neg j)(j)
\]
both inhabit $\Path\,(\Path\,A\,a_{10}\,a_{11})\,m\,m$, together
with the edge square
$S_m := \langle i \rangle \langle j \rangle\, m(j)$, and are
pairwise non-convertible \cite{Realization}.

\begin{prop}[No direct homotopy]\label{prop:nodirect}
  There is no strict homotopy with $k$-faces $S_1$ and $S_2$: by
  the realization theorem a strict cube is a single cell's
  material, and its two $k$-faces are that cell's material or
  boundary material --- never the interiors of two different
  cells.
\end{prop}

\begin{thm}[Connection through the shared subcomplex]
  \label{thm:zigzag}
  Both interiors deform strictly onto the shared material:
  \[
    \langle k \rangle \langle i \rangle \langle j \rangle\,
    Q_1\bigl((i \vee \neg i \vee j \vee \neg j) \vee k\bigr)(j)
    \;:\; S_1 \rightsquigarrow S_m,
  \]
  \[
    \langle k \rangle \langle i \rangle \langle j \rangle\,
    Q_2\bigl((i \wedge \neg i \wedge j \wedge \neg j) \wedge
    \neg k\bigr)(j)
    \;:\; S_2 \rightsquigarrow S_m,
  \]
  --- at $k = 1$ the first formula becomes $1$ and
  $Q_1(1)(j) \aconv m(j)$ by the attachment equations of
  \cite{Realization}; dually for the second.  Hence the strict
  fiber is connected by the zig-zag
  $S_1 \rightsquigarrow S_m$, $S_2 \rightsquigarrow S_m$ --- entirely inside
  the strict layer.  (Machine witnesses \texttt{mixDeformQ1D},
  \texttt{mixDeformQ2D}.)
\end{thm}

\begin{rem}[The general pattern]\label{rem:descent}
  When same-boundary fillers originate from different cells, the
  realization theorem forces the boundary to lie in the
  \emph{shared} subcomplex, and the coordinate of each filler
  transverse to its shared face then has endpoint restrictions on
  every face --- so Lemma~\ref{lem:interp} applies against the
  corresponding constant, deforming each filler onto shared
  material, where Theorem~\ref{thm:contract} (one dimension down)
  connects them.  We have verified the two-cell, dimension-$2$
  case by machine; the inductive general statement follows the
  same descent and we record it at the same level of uniformity as
  the realization theorem itself \cite[Rem.~5.4]{Realization}.
\end{rem}

\section{The strict--weak comparison}
\label{sec:comparison}

The full theory adds the Kan operations.  Two theorems close the
comparison over a common boundary.

\begin{thm}[$(\dagger)$ is propositional]\label{thm:dagger}
  Let $T := \langle j \rangle\,
  \hcomp\,[\,j{=}0 \mapsto a,\ j{=}1 \mapsto b\,]\,(p\,j)$ be the
  constant-tube composite over a generic edge $p$ --- the left-hand
  side of the equation $(\dagger)$ that the no-go theorem of
  \cite{Const} proves can never be definitional.  Then the
  canonical filler
  \[
    \langle k \rangle \langle j \rangle\,
    \hcomp\,[\,j{=}0 \mapsto a,\ j{=}1 \mapsto b,\
    k{=}0 \mapsto p\,j\,]\,(p\,j)
    \;:\; \Path\,(\Path\,A\,a\,b)\,p\,T
  \]
  is a path rel boundary from $p$ to $T$: at $k = 0$ the added face
  is decided and the composite collapses to $p\,j$; at $k = 1$ the
  false face is discarded, leaving $T$.  The same two computations
  give the statement for any constant-tube system in any dimension.
  (Machine witness \texttt{mixFillD}, together with the standing
  definitional separation $T \nconv p$.)
\end{thm}

\begin{rem}
  Theorem~\ref{thm:dagger} is the exact dual of the no-go theorem:
  \emph{what can never be definitional is precisely what
  $\hfillop$ supplies propositionally, uniformly}.  The
  definitional/propositional boundary for composition is not a
  deficiency to be engineered away --- the no-go theorem says it
  cannot be --- but a structural line, with the canonical filler
  standing exactly on it.
\end{rem}

\begin{prop}[The weak layer extends the propositional
  classification]\label{prop:winding}
  In the circle instance, the composite
  $\mathsf{loop} \cdot \mathsf{loop}$ has winding number $+2$,
  while $\refl$, $\mathsf{loop}$ and its reversal have winding
  $0, +1, -1$ (all machine-computed in the artifact library of
  \cite{Const}).  Since winding is a propositional invariant,
  genuine composites are not propositionally equal to these strict
  cells: the weak layer adds new propositional classes over the
  same boundary.
\end{prop}

Table~\ref{tab:comparison} now stands proved in all three rows.

\section{Machine verification}
\label{sec:artifact}

All instances are elaboration-time checks in the self-contained
Lean~4 cubical kernel of \cite{Const}: repository
\begin{center}
  \url{https://github.com/r-shibusawa/cubical-strict-j}.
\end{center}
The probes of this paper are the file \texttt{LibMixed.lean} ---
the interpolation homotopy with its control and separation, the
machine-checked \emph{failure} of the naive interpolant, the two
deformation homotopies of the zig-zag, and the canonical filler
with its definitional-separation control --- together with the
proof notes \texttt{docs/DeMorganConvexity.md} and
\texttt{docs/MixedLayer.md}.  They are included from release
\texttt{v1.4.0}, archived at DOI
\href{https://doi.org/10.5281/zenodo.TODOE}%
{\texttt{10.5281/zenodo.\allowbreak TODOE}}.  A successful
\texttt{lake build} replays every check; by (K3) each check is a
decision.

\section{Related work and outlook}
\label{sec:related}

The uniqueness of Kan fillers up to homotopy is foundational for
cubical models \cite{CCHM}; the present paper locates its exact
strict shadow, running on the connection algebra alone.  The
series' platform and characterizations are
\cite{Const,Sphere,Realization}; winding numbers and the circle
computation descend from the encode--decode method
\cite{LicataShulman}.

Two directions seem natural.  First, \emph{a strict model
structure in miniature}: De Morgan convexity gives the strict
layer contractible fibers and canonical homotopies; whether this
organizes into a factorization system or a fibration structure on
cellular presheaves --- with the strict layer as the ``cofibrant
skeleton'' of the full theory --- deserves study.  Second,
\emph{normal forms up to homotopy}: combining
Theorem~\ref{thm:dagger} with the realization theorem suggests
rewriting genuinely weak cells to canonical composites of strict
cells, a step toward word problems for cubical polygraphs
\cite[\S 10]{Realization}.

\subparagraph*{Acknowledgements.}  The platform and companion
results were developed in \cite{Const,Sphere,Realization}.

\begin{thebibliography}{10}

\bibitem{CCHM}
C.~Cohen, T.~Coquand, S.~Huber, A.~M\"ortberg.
\newblock Cubical type theory: a constructive interpretation of the
  univalence axiom.
\newblock \emph{TYPES 2015}, LIPIcs 69, 2018.

\bibitem{LicataShulman}
D.~R.~Licata, M.~Shulman.
\newblock Calculating the fundamental group of the circle in
  homotopy type theory.
\newblock \emph{LICS 2013}.

\bibitem{Const}
R.~Shibusawa.
\newblock Evaluation-time constancy inference for transport in
  cubical type theory.
\newblock Preprint, 2026.  Artifact:
  \url{https://github.com/r-shibusawa/cubical-strict-j}, archived
  at DOI \href{https://doi.org/10.5281/zenodo.21637898}%
  {10.5281/zenodo.21637898}.

\bibitem{Sphere}
R.~Shibusawa.
\newblock The strict layer of generic cells in cubical type theory
  is a wedge of De Morgan spheres.
\newblock Preprint, 2026.  Included in the artifact repository
  above (from release \texttt{v1.3.0}, DOI
  \href{https://doi.org/10.5281/zenodo.21638976}%
  {10.5281/zenodo.21638976}).

\bibitem{Realization}
R.~Shibusawa.
\newblock Generic cellular contexts in cubical type theory realize
  their cube presheaves on the nose.
\newblock Preprint, 2026.  Included in the artifact repository
  above (from release \texttt{v1.3.0}, DOI as for \cite{Sphere}).

\end{thebibliography}

\end{document}
